\subsection{Exact affine and line lift comparison at additive caps}

\begin{statement}{P3F.44}{Proved}
Every actual affine line lift on one of the constructed additive root sources is carried by the explicit source translation and line maps below to $\zeta_m^d A_0$, with $d\bmod m$ retained. The comparison preserves the actual suspension $B_0$ and transports the outer annular map. It does not discard the transported source translation or integral line gauge.
\end{statement}
\begin{proof}
\par

\medskip\noindent\textbf{AE.1. The section group and its exact cyclic action}

Let \(D\) be the canonical order-\(m\) root deck on the
constructed minimal regular source \(S'\to\Delta_z\). Its
action on an analytic section is
\[
 (D\cdot T)(z)=D\bigl(T(\zeta_m^{-1}z)\bigr).
\]
All sums below are sums in this actual section group. The marked
section \(Q'\) is \(D\)-invariant. The constructed branch
\(R\), with \(mR=Q'\), is also \(D\)-invariant. At a
reduced II, III or IV root it is the unique division branch
because the central identity group is additive and multiplication
by \(m\) is invertible on the component group. At a good root
it is the branch with central value \(O\).

Every affine source map with the specified linear deck is
\(a=t_S D\) for an actual analytic section \(S\). Its exact
norm equation is
\[
 a^m=t_{-nQ'},\qquad
 \sum_{j=0}^{m-1}D^jS=-nQ'.                         \tag{AE.1}
\]
Set \(Z=S+nR\). Then \(\sum D^jZ=0\). Conjugation of the
constructed source map \(a_0=t_{-nR}D\) by translation is
\[
 t_T a_0t_T^{-1}=t_{T-DT-nR}D.                    \tag{AE.2}
\]
Thus the precise required equation is \((1-D)T=Z\).

\par

\medskip\noindent\textbf{AE.2. An explicit solution on every constructed root}

At a reduced root multiplication by \(m\) is a bijection on
germs of analytic sections. On the additive central identity
group its inverse is multiplication by \(1/m\), and on the
finite component group it is a bijection because \(m\) is
prime to its order. Each central inverse extends uniquely by
the analytic inverse function theorem. An analytic section in
its kernel has central value zero, so that same uniqueness
shows that it is the zero section. Define
\[
 T=-[m]^{-1}\left(\sum_{j=0}^{m-1}jD^jZ\right).
                                                               \tag{AE.3}
\]
Since \(D^m=1\), direct subtraction gives
\[
 (1-D)\sum_{j=0}^{m-1}jD^jZ
   =\sum_{j=0}^{m-1}D^jZ-mZ=-mZ.
\]
Equation (AE.3) therefore solves (AE.2) exactly. The translation
extends over the resolved minimal regular source by its actual
source group law and graph comparison, already used to extend
the norm map in SC/OA/II.

At a good root the central linear automorphism \(D_0\) has
order three, four or six and has no eigenvalue one. Hence
\(1-D_0\) is a surjective elliptic isogeny. Choose a central
point \(T_0\) with
\((1-D_0)T_0=Z_0\), and lift it to an analytic section
\(\widetilde T\) of the smooth elliptic family. Put
\[
 Z_1=Z-(1-D)\widetilde T.
\]
It has central value zero and still has zero cyclic norm.
The subgroup of analytic sections with central value zero has
unique \(m\)-division with central value zero, by the analytic
inverse function theorem at \(O\). Apply (AE.3) to \(Z_1\)
in this subgroup and call the result \(T_1\). Then
\(T=\widetilde T+T_1\) solves \((1-D)T=Z\). Substitution
in (AE.2) is the required actual source isomorphism. No choice
of a rational homology marking was used to obtain it.

This also verifies the good-root norm restriction without losing
the central point. On the good central fibre
\(1+D_0+\cdots+D_0^{m-1}=0\), so (AE.1) forces
\(nQ'_0=0\). For III/IV, \(Q'_0\) has order two/three
when its component is nonzero; coprimality with \(m\) excludes
that case. For II its actual order-six fixed-point calculation
gives \(Q'_0=O\), and the constructed norm is admissible.

\par

\medskip\noindent\textbf{AE.3. The actual translation lift on the specified line}

Let \(N\) be the resolved root line already constructed. On
every reduced central component it has degree zero, and on a
good central fibre it has degree zero. Translation by \(T\)
preserves these degree vectors. On a smooth generic fibre,
\(t_T^*N\otimes N^{-1}\) is trivial: in a degree-zero
uniformization its factors of automorphy are constant in the
elliptic variable, so a lift of the translation leaves them
unchanged.

There is an extending holomorphic line isomorphism on the whole
source, not only a generic Picard-class equality. Apply proper
direct-image coherence to
\(L=t_T^*N\otimes N^{-1}\). Its direct image is torsion-free
of generic rank one on the disc, hence free after shrinking.
A nonzero generator evaluates to a section of \(L\) whose
divisor is vertical: on each generic compact elliptic fibre it
is a nonzero section of a trivial degree-zero line. Its degree
on each central component is zero. The exact component
intersection matrix consequently puts that divisor in its
primitive whole-fibre kernel. At these reduced or good roots
the fibre has multiplicity one, so dividing by its integral
base monomial removes the divisor. The resulting section is
nowhere zero and gives the required line map \(\widetilde t_T\).

This lift commutes with the actual suspension \(B_0\). On a
generic elliptic fibre write the degree-zero line as the quotient
of \(\mathbb C\times\mathbb C^*\) by constant factors of
automorphy. Source translations lift as
\((z,v)\mapsto(z+t,cv)\); these maps commute. The commutator
of the extending maps is a holomorphic unit on the root source.
It is one on every generic fibre and hence one everywhere.
Thus conjugating by \(\widetilde t_T\) transports the actual
line root while preserving \(B_0\), not merely its source
translation. The base monomial used above retains its deck
character in this transported map.

\par

\medskip\noindent\textbf{AE.4. The remaining finite delta character and the boundary map}

After this source-and-line conjugation, compare the given line
root \(A\) with the constructed \(A_0\), which have the same
source map. Define the ratio by \(AA_0^{-1}=M_c\), where
\(M_c\) multiplies the line at base coordinate \(z\) by
\(c(z)\). This is a nowhere-zero holomorphic function
on \(S'\), hence a base unit \(c(z)\), because the source
fibres are compact and connected. The equality of their exact
norms is
\[
 \prod_{j=0}^{m-1}c(\zeta_m^jz)=1.                \tag{AE.4}
\]
Its central value is \(c(0)=\zeta_m^d\) for a retained
integer \(d\bmod m\). Divide this constant and take the
analytic logarithm \(h(z)\) with \(h(0)=0\). Equation
(AE.4) implies \(\sum D^jh=0\); the possible integral
logarithm constant is zero at the centre. The additive formula
\[
 b(z)=-\frac1m\sum_{j=0}^{m-1}jD^jh(z)
\]
gives \((1-D)b=h\), with \((Db)(z)=b(\zeta_m^{-1}z)\).
The complete conjugation identity, with its arrow fixed, is
\[
 A=M_{e^b}\,(\zeta_m^dA_0)\,M_{e^{-b}},
\]
because its scalar multiplier is
\(\zeta_m^d\exp(b(z)-b(\zeta_m^{-1}z))=c(z)\).
It leaves \(B_0\) unchanged
because that map fixes the base.

Consequently every actual affine line lift on each constructed
root is carried to \(\zeta_m^dA_0\) by the displayed maps,
with its finite delta character retained. Their restrictions to
the outer annulus are the boundary comparison maps. They may
contain a source translation and an integral line gauge; these
are transported, not set to zero. The one-cap topology proofs
allow the resulting integer collar coefficient in full, and
their degree-two matrix remains unimodular for every value.
Global parameter-return comparisons continue to use the actual
closing map and these transported characters, as in K24.
\end{proof}
